\documentclass[11pt]{article}

\usepackage[margin=1in]{geometry}
\usepackage{amsmath,amssymb,amsthm,mathtools}
\usepackage{booktabs}
\usepackage{enumitem}
\usepackage{xcolor}
\usepackage{hyperref}
\usepackage[T1]{fontenc}

\hypersetup{colorlinks=true,linkcolor=blue,citecolor=blue,urlcolor=blue}

\newtheorem{theorem}{Theorem}[section]
\newtheorem{proposition}[theorem]{Proposition}
\newtheorem{lemma}[theorem]{Lemma}
\newtheorem{corollary}[theorem]{Corollary}
\theoremstyle{definition}
\newtheorem{definition}[theorem]{Definition}
\newtheorem{remark}[theorem]{Remark}

\newcommand{\Z}{\mathbb Z}
\newcommand{\F}{\mathcal F}
\newcommand{\kg}{\mathrm{KG}}
\newcommand{\MSF}{\mathrm{MSF}}
\DeclareMathOperator{\LCM}{lcm}

\title{Erd\H{o}s Problem \#7: Odd Distinct Covering Systems\\
\large BFF reconstruction, finite certificates, and prime-power refinements}
\author{Research working draft}
\date{12 August 2026}

\begin{document}
\maketitle

\begin{center}
\fcolorbox{red!60!black}{red!5}{\parbox{0.92\textwidth}{
\textbf{Status warning.} This is a working research document, not a paper or novelty claim.
Statements explicitly labelled ``published'' are traced to cited sources.  Statements labelled
``derived'' are deductions made in this project and should be independently checked before
circulation.  The finite computations are reproduced by the accompanying exact-rational Python
scripts.  In particular, the bound $218{,}243{,}025$ and the maximum-spanning-forest refinement
have not been peer reviewed and are not asserted here to be absent from the prior literature.
}}
\end{center}

\begin{abstract}
Erd\H{o}s and Selfridge asked whether there exists a finite covering system of congruences
$a_i\pmod{m_i}$ whose moduli are pairwise distinct, greater than one, and all odd.  We organize
the problem in finite cyclic-group, product-space, divisor-lattice, hypergraph, probabilistic,
and optimization languages.  We then reconstruct the exponent-sensitive necessary inequality
of Berger, Felzenbaum and Fraenkel (BFF), emphasizing that their correction term is a forest
bound on forced intersections.  This immediately suggests a derived maximum-spanning-forest
refinement.  Combining published structural restrictions with exact enumeration of exponent
vectors gives a provisional, fully reproducible deduction that a hypothetical odd covering has
LCM at least $218{,}243{,}025$.  The boundary shape
$3^3 5^2 7\cdot11\cdot13\cdot17\cdot19$ survives both the published BFF inequality and the
first forest/triangle refinements, identifying a concrete target for further work.  A block or
``tower capacity'' lemma is introduced as a second route for exploiting prime-power depth.
\end{abstract}

\tableofcontents

\section{The problem and finite reformulations}

A \emph{covering system} is a finite family
\[
 a_i\pmod{m_i},\qquad 1\le i\le k,
\]
whose union is all of $\Z$.  It is \emph{distinct} when the $m_i>1$ are pairwise distinct.
Erd\H{o}s Problem \#7 asks whether a distinct covering can have every $m_i$ odd.

Let
\[
 L=\LCM(m_1,\dots,m_k).
\]
The covering is periodic modulo $L$, so it is equivalent to a cover of $\Z/L\Z$ by cosets of
subgroups.  If
\[
 L=\prod_{i=1}^{n}p_i^{s_i},\qquad p_1<\cdots<p_n,
\]
then the Chinese remainder theorem gives
\[
 \Z/L\Z\cong\prod_{i=1}^{n}\Z/p_i^{s_i}\Z.
\]
A divisor $m=\prod p_i^{d_i}$ is naturally encoded by a depth vector
$\mathbf d=(d_1,\dots,d_n)$, $0\le d_i\le s_i$.  Thus a residue class modulo $m$ is a cylinder
fixing $d_i$ $p_i$-adic digits in coordinate $i$.  The square-free case is the special case
$s_i=1$ for every $i$; Balister et al. proved that an odd distinct covering is impossible in
that case \cite{BalisterSquarefree2021}.  Consequently the unresolved problem is intrinsically
multilevel or prime-power in nature.

The same finite object can be viewed as a constrained set-cover problem, a hypergraph cover,
or a probability problem on uniformly random $X\in\Z/L\Z$.  For the event
$A_i=\{X\equiv a_i\pmod{m_i}\}$, $\Pr(A_i)=1/m_i$, while the Chinese remainder theorem controls
many intersections.  These translations motivate both the BFF overlap inequality and modern
SAT/MIP approaches.

\section{Published structural restrictions}

We record only the restrictions used later.

\begin{theorem}[BFF II, published \cite{BFF1987}]
If an odd incongruent covering system exists and the LCM has $n$ distinct prime factors, then
$n\ge6$.  More precisely, the exponent-sensitive parameters defined in Section~\ref{sec:bff}
must satisfy $g(\bar w,\bar{\mathbf z})\ge2$.
\end{theorem}

\begin{theorem}[Hough--Nielsen, published \cite{HoughNielsen2019}]
Every distinct covering system contains a modulus divisible by $2$ or $3$.
Hence, for an odd distinct covering,
\[
 3\mid L.
\]
\end{theorem}

\begin{theorem}[Balister et al., published \cite{BalisterSquarefree2021}]
There is no distinct covering system with all moduli odd and square-free.  In particular, the
LCM of a hypothetical odd distinct covering cannot be square-free.
\end{theorem}

A different theorem of Balister et al. gives the useful restriction that an odd covering LCM
must be divisible by $9$ or $15$ \cite{BalisterDensity2022}.  It is not required for the
finite bound derived below, but it is consistent with the first surviving shape.

\section{Re-deriving the BFF inequality}\label{sec:bff}

\subsection{The product box and building blocks}

Write
\[
 L=\prod_{i=1}^{n}p_i^{s_i},\qquad P_i=p_i^{s_i}.
\]
BFF represent the relevant finite group by the product box
\[
 \mathcal P=[0,P_1)\times\cdots\times[0,P_n),
 \qquad |\mathcal P|=L,
\]
and refine cosets into product ``building blocks'' indexed by nonempty subsets
$I\subseteq[n]$.  The index records coordinates on which the block is restricted.
After enlarging the family in a way that can only make covering easier, the number of
blocks of a given index is bounded by BFF's quantity $\alpha(I)$.  The exceptional counts
for one-coordinate indices are the source of the asymmetric first coordinate in the final
formula \cite{BFF1987}.

Remove the blocks whose index has cardinality one.  The remaining set $\mathcal S$ is again a
product set.  Its coordinate sizes are
\[
 y_1=P_1-\frac{P_1-1}{p_1-1}
     =\frac{(p_1-2)P_1+1}{p_1-1},
\]
and, for $i\ge2$,
\[
 y_i=P_i-2\frac{P_i-1}{p_i-1}
     =\frac{(p_i-3)P_i+2}{p_i-1}.
\]
This normalization makes the BFF parameters transparent:
\begin{align}
 w&=\frac{P_1-1}{(p_1-2)P_1+1},\\
 z_1&=\frac{p_1^{s_1-1}-1}{(p_1-2)P_1+1},\\
 z_i&=\frac{P_i-1}{(p_i-3)P_i+2},\qquad 2\le i\le n.
\end{align}
They are normalized capacities of higher-index blocks inside the surviving box.

\subsection{The pre-forest union bound}

For $|I|\ge2$, BFF bound the fraction of $\mathcal S$ occupied by all blocks of index $I$.
Summing these bounds and collecting terms gives
\begin{equation}\label{eq:g1}
 g_1(w,\mathbf z)
 =(1+w)\prod_{i=2}^{n}(1+z_i)-w
 -(1+w-z_1)\sum_{i=2}^{n}z_i.
\end{equation}
The naive argument would require $g_1\ge2$: the higher-index blocks must cover the entire
residual product set, and $g_1-1$ is their normalized union-bound capacity.

\subsection{Forced intersections and the forest lemma}

Let $\mathcal J=\binom{[n]}2$.  For $I,J\in\mathcal J$ with $I\cap J=\varnothing$, BFF show
under their worst-case enlargement that the corresponding type-$I$ and type-$J$ unions have
forced intersection of normalized size
\begin{equation}\label{eq:beta}
 \beta(I,J)=\prod_{r\in I\cup J}z_r.
\end{equation}
They use the elementary forest inequality
\[
 \left|\bigcup_{v\in V}S_v\right|
 \le \sum_{v\in V}|S_v|-
 \sum_{\{u,v\}\in E(T)}|S_u\cap S_v|
\]
for any forest $T$ on the set family.  Therefore every forest on $\mathcal J$ whose edges join
disjoint 2-subsets gives a legitimate forced-overlap correction.

BFF choose one particular forest supported on the first five coordinates, with total weight
\begin{equation}\label{eq:wbff}
 W_{\rm BFF}
 =3z_1z_2z_3z_4+3z_1z_2z_3z_5
  +2z_1z_2z_4z_5+z_1z_3z_4z_5.
\end{equation}
Thus their published polynomial is
\begin{equation}\label{eq:gbff}
 g_{\rm BFF}=g_1-W_{\rm BFF},
\end{equation}
and a necessary condition for an odd distinct cover is
\begin{equation}\label{eq:bffcondition}
 g_{\rm BFF}\ge2.
\end{equation}
This re-derivation highlights the conceptual content: BFF is a naive capacity estimate minus a
rigorously selected set of unavoidable overlaps.

\section{Derived maximum-spanning-forest refinement}

The published forest lemma does not require BFF's particular forest.  This yields an immediate
optimization.

\begin{definition}
Let $G_n$ be the Kneser disjointness graph $\kg(n,2)$: vertices are the 2-subsets of $[n]$ and
$I$ is adjacent to $J$ precisely when $I\cap J=\varnothing$.  Give each edge the weight
$\beta(I,J)$ from \eqref{eq:beta}.  Let $W_{\MSF}(\mathbf z)$ be the maximum total weight of a
forest in $G_n$.
\end{definition}

\begin{proposition}[Derived]
Every odd distinct covering satisfies
\begin{equation}\label{eq:gmsf}
 g_{\MSF}(w,\mathbf z)
 :=g_1(w,\mathbf z)-W_{\MSF}(\mathbf z)\ge2.
\end{equation}
Moreover $g_{\MSF}\le g_{\rm BFF}$.
\end{proposition}

\begin{proof}
BFF's forest lemma and their exact intersection formula apply to every forest with edges between
disjoint 2-index sets.  Choose a maximum-weight such forest.  The BFF forest is one feasible
choice, hence $W_{\MSF}\ge W_{\rm BFF}$ and the resulting necessary condition is no weaker.
\end{proof}

For $n\ge5$, $\kg(n,2)$ is connected, so with positive weights the optimum can be taken as a
maximum-weight spanning tree.  The accompanying script `msf\_refinement.py` constructs it by
exact-rational Kruskal optimization.

As a minimal hand-check that BFF's displayed tree need not be optimal, one may replace the last
term of \eqref{eq:wbff} by $z_2z_3z_4z_5$ using a different admissible tree.  The gain is
$(z_2-z_1)z_3z_4z_5$ whenever $z_2>z_1$; otherwise retain the published tree.

\subsection{Connection with Hunter--Worsley bounds: where the pairwise route stops}

The forest inequality used by BFF is the same structural mechanism as the Hunter--Worsley
second-order upper bound for a union: subtract a maximum-weight spanning tree of pairwise
intersections from the sum of the singleton probabilities \cite{Hunter1976,Worsley1982}.
Modern probability-optimization literature sharpens the interpretation.  For fixed singleton
probabilities together with only lower bounds on pairwise intersections, the optimal upper union
bound is of Hunter type; Ramachandra and Natarajan summarize this result and its relation to
Boros et al. \cite{BorosEtAl2014,RamachandraNatarajan2023}.

Our BFF setting is not literally the generic fixed-marginal probability problem---the singleton
quantities in \eqref{eq:g1} arise from an arithmetic relaxation---so we do not import an
``optimality theorem'' blindly.  Nevertheless, after that relaxation, the only additional data
used by the forest step are lower bounds on pair intersections for disjoint 2-index groups.  The
maximum-spanning-forest refinement therefore exhausts the obvious Hunter--Worsley improvement of
that pairwise data.  This is an important negative conclusion: a substantial next improvement
should use information that BFF's pairwise forest discards, such as exact triple intersections,
compatibility among coordinate projections, higher-index groups, or the prime-power fibre
structure.

\section{A provisional exact lower bound for the LCM}\label{sec:218m}

This section combines published theorems with a finite exact computation.  It is deliberately
labelled provisional pending independent verification and a more exhaustive novelty search.

\subsection{Prime monotonicity for fixed exponent pattern}

For $i\ge2$ define
\[
 z(p,s)=\frac{p^s-1}{(p-3)p^s+2}.
\]
Treating $p$ temporarily as a real variable,
\begin{equation}\label{eq:zderivative}
 \frac{\partial z}{\partial p}
 =\frac{p^{s-1}\left(s(p-1)+p-p^{s+1}\right)}{\left((p-3)p^s+2\right)^2}.
\end{equation}
For $p\ge5$ and $s\ge1$, the numerator is negative, so $z(p,s)$ decreases with $p$.
BFF state that $g_{\rm BFF}$ is increasing in its variables on the relevant domain
\cite{BFF1987}.  Since Hough--Nielsen forces $p_1=3$, replacing any later ordered prime by a
smaller allowed prime can only increase $g_{\rm BFF}$ for a fixed exponent vector.  The
$s_1=1$ boundary, for which $z_1=0$, follows by evaluating the same polynomial on the closure of
the domain.

\subsection{Finite exponent enumeration}

Set
\begin{equation}\label{eq:X}
 X=3^3 5^2 7\cdot11\cdot13\cdot17\cdot19
  =218{,}243{,}025.
\end{equation}

For exactly six distinct prime factors, it suffices by the preceding monotonicity to enumerate
positive exponent vectors for the minimal tuple
\[
 (3,5,7,11,13,17)
\]
whose minimal-prime shape is below $X$.  There are 89 such vectors.  Exact rational evaluation
of \eqref{eq:gbff} gives $g_{\rm BFF}<2$ in every case.  The largest value is
\[
 \frac{506755}{255717}\approx1.98170242886
\]
at exponent vector $(4,3,1,1,1,1)$, whose minimal-prime shape is $172{,}297{,}125$.

For exactly seven distinct prime factors, the minimal tuple is
\[
 (3,5,7,11,13,17,19).
\]
There are 16 positive exponent vectors below $X$.  Again every one has $g_{\rm BFF}<2$; the
maximum is
\[
 \frac{117043}{60775}\approx1.92584121761
\]
at $(2,2,1,1,1,1,1)$, with shape $72{,}747{,}675$.

For eight or more distinct primes, nonsquarefreeness and $3\mid L$ imply
\[
 L\ge3^2\cdot5\cdot7\cdot11\cdot13\cdot17\cdot19\cdot23
   =334{,}639{,}305>X.
\]

\begin{proposition}[Derived + exact finite certificate; provisional]
Assuming the published BFF, Hough--Nielsen, and square-free theorems as quoted above, the finite
calculation in `verify\_218m\_bound.py` yields
\[
 \boxed{L\ge218{,}243{,}025}
\]
for the LCM of any hypothetical odd distinct covering system.
\end{proposition}

\begin{proof}
BFF gives at least six distinct prime factors; Hough--Nielsen gives $3\mid L$; and the
square-free theorem forces at least one exponent to exceed one.  If $L<X$, then $n\ge8$ is
impossible by the displayed numerical lower bound.  For $n=6$ and $n=7$, prime monotonicity
reduces every fixed exponent pattern to the corresponding minimal prime tuple, and the complete
positive-exponent enumerations below $X$ have respectively 89 and 16 cases.  The accompanying
CSV records the exact fraction $g_{\rm BFF}$ for every case; all are strictly below two,
contradicting the BFF necessary condition.
\end{proof}

The script uses Python's `fractions.Fraction`; floating-point numbers are printed only for human
readability.  The CSV is therefore an explicit finite arithmetic certificate, though not a
formal proof-assistant certificate.

\section{The first surviving boundary shape}

At the boundary
\begin{equation}\label{eq:L0}
 L_0=3^3 5^2 7\cdot11\cdot13\cdot17\cdot19
     =218{,}243{,}025,
\end{equation}
the BFF parameters are
\[
 w=\frac{13}{14},\qquad
 (z_1,\dots,z_7)=
 \left(\frac27,\frac6{13},\frac15,\frac19,\frac1{11},\frac1{15},\frac1{17}\right).
\]
Exact evaluation gives
\[
 g_{\rm BFF}=\frac{175089}{85085}\approx2.05781277546>2,
\]
so the published inequality no longer excludes the shape.

The maximum-spanning-forest refinement gives
\[
 g_{\MSF}=\frac{2608283}{1276275}\approx2.04366848837>2.
\]
Thus forest optimization closes about a quarter of the BFF excess above two, but not enough to
eliminate $L_0$.

\section{Going one order beyond forests: an exact triangle cluster}

When $n\ge6$, $\kg(n,2)$ contains triangles: three pairwise-disjoint 2-subsets
$I,J,K$.  For such a triple, product structure fixes not only all three pair intersections but
also the triple intersection.  Hence the union of these three type groups admits exact
third-order inclusion--exclusion,
\[
 |A_I\cup A_J\cup A_K|
 =\sum |A|-\sum |A\cap A'|+|A_I\cap A_J\cap A_K|,
\]
with normalized triple intersection
\[
 \gamma(I,J,K)=\prod_{r\in I\cup J\cup K}z_r.
\]
Contract the triangle to one cluster.  The remaining groups can again be attached by a forest;
for an edge from the cluster to an outside group, use any one triangle member whose guaranteed
intersection lies inside the cluster intersection, and choose the largest such lower bound.
This gives a rigorous derived refinement.

At $L_0$, exhaustive exact search over triangle choices gives
\[
 g_{\triangle}=\frac{23472541}{11486475}
 \approx2.04349384820>2.
\]
The gain beyond the maximum spanning forest is real but small.  This experiment indicates that
higher-order type-2 inclusion--exclusion by itself is unlikely to close the remaining gap; the
prime-power depth should enter more directly.

\section{A block or tower capacity lemma}

McNew and Setty define $r(n)$ as the largest number of residues modulo $n$ coverable by distinct
nontrivial divisors of $n$, and $c(n)=1+r(n)/n$.  Thus $n$ is a covering number exactly when
$c(n)=2$.  They prove
\[
 c(mn)\le c(m)h(n),\qquad (m,n)=1,
 \qquad h(n)=\frac{\sigma(n)}n,
\]
and conjecture the stronger coprime submultiplicativity $c(mn)\le c(m)c(n)$
\cite{McNewSetty2026}.  The following elementary block lemma is tailored to the CRT product.

\begin{lemma}[Derived block-capacity certificate]\label{lem:block}
Let
\[
 N=M_1\cdots M_t,
\]
where the $M_j$ are pairwise coprime.  Suppose $R_j$ is an upper bound for $r(M_j)$.  Call a
nontrivial divisor $d\mid N$ \emph{pure} if $d\mid M_j$ for some $j$, and \emph{mixed}
otherwise.  If
\begin{equation}\label{eq:blockcertificate}
 \sum_{\substack{d\mid N\\d\text{ mixed}}}\frac1d
 <
 \prod_{j=1}^{t}\left(1-\frac{R_j}{M_j}\right),
\end{equation}
then $N$ is not a covering number.
\end{lemma}

\begin{proof}
Pure congruences supported inside block $M_j$ cover at most $R_j$ residues in that coordinate.
By CRT, after all pure congruences there remains a Cartesian product of coordinate complements
of relative size at least the right-hand side of \eqref{eq:blockcertificate}.  A mixed modulus
$d$ covers exactly a fraction $1/d$ of the full space, so all mixed moduli together cover at
most the left-hand side by the union bound.  If that capacity is smaller than the guaranteed
pure-block residual, some point remains uncovered.
\end{proof}

The left side is easy to compute:
\[
 \sum_{d\text{ mixed}}\frac1d
 =h(N)-1-\sum_{j=1}^{t}\bigl(h(M_j)-1\bigr).
\]
For prime powers one has the exact formula
\[
 r(p^a)=1+p+\cdots+p^{a-1}=\frac{p^a-1}{p-1},
\]
since the union bound is attained by choosing pairwise-disjoint residue classes at successive
$p$-adic depths.

At $L_0$, use the tower blocks
\[
 27,25,7,11,13,17,19.
\]
The guaranteed pure-block residual is
\[
 \frac{3072}{12155}\approx0.25273549979,
\]
whereas the mixed reciprocal budget is
\[
 \frac{6954221}{11486475}\approx0.60542690425.
\]
Thus the simplest block certificate fails.  This failure is informative: a successful
prime-power argument must either absorb many mixed moduli into larger blocks, improve the
block capacities $r(M)$ substantially, or prove significant unavoidable overlap among the
mixed moduli.

\section{A diagnostic for the 3-adic fibre route}

The preceding failure suggests refining by the $3$-adic depth.  Before building a large model,
it is useful to determine whether mere reciprocal-capacity bookkeeping across fibres could ever
be enough.  Write
\[
 L_0=27Q,\qquad Q=25\cdot7\cdot11\cdot13\cdot17\cdot19=8{,}083{,}075.
\]
A modulus $3^a d$ with $(d,3)=1$ is compatible with $27/3^a$ fibres modulo $27$ and occupies a
fraction $1/d$ inside each compatible fibre.  At depth $a=0$, the union bound over all
nontrivial $d\mid Q$ gives the uniform fibre capacity
\[
 C_0=\sum_{\substack{d\mid Q\\d>1}}\frac1d
 =\frac{197599}{230945}\approx0.85561064323.
\]
Hence a capacity-only argument leaves deficit $1-C_0\approx0.14438935677$ in each fibre.
At depth one, assign the allowed cofactor divisors $\{1\}$, $\{5\}$, and $\{7,11\}$ to the
three residue branches modulo $3$.  These contribute respectively
\[
 1,\qquad \frac15,\qquad \frac17+\frac1{11},
\]
each exceeding the deficit in all nine descendant fibres of its branch.  Therefore the relaxed
fibre-capacity model already permits normalized load at least one in every fibre without even
using depths two or three.

This is a useful negative result: a serious $3$-adic argument must retain the cofactor coverage
sets, or at least strong certified bounds on their simultaneous realizability.  Reciprocal sums
per fibre alone cannot exclude $L_0$.  The calculation is reproduced by
\texttt{fiber\_relaxation.py}.

\section{Relation to the current computational frontier}

Mian and Siddique formalized finite exclusions in Lean and proved, with a kernel-checked
certificate chain, that any odd covering LCM exceeds $10{,}000$ \cite{MianSiddique2026}.  Their
paper emphasizes a finite CRT bridge designed to consume future search certificates.  McNew and
Setty use optimization to push the classification of primitive covering numbers essentially
through $10^6$, with one unresolved even case in that range \cite{McNewSetty2026}.  These works
and the present BFF-based calculation have different trust models and purposes: the aim here is
a compact structural certificate rather than a large black-box classification.

Recent constructive work also shows that the odd-modulus problem is close to solvable when one
permits limited repetition; for example, constructions with only a small number of repeats of a
chosen odd modulus are studied in \cite{Bispels2025}.  Such constructions reinforce the view
that prime-power branching is central rather than incidental.

\section{Next proof targets}

The numerical checkpoint at $L_0$ gives a focused research program.

\begin{enumerate}[label=\arabic*.]
\item \textbf{3-adic fibre inequality.}  Split $\Z/L_0\Z$ into its 27 fibres modulo $27$ and
track moduli of 3-adic depth $0,1,2,3$.  Distinctness restricts how much cofactor capacity can be
concentrated in the same fibres.  A successful inequality here would use information discarded
by the scalar BFF parameters.

\item \textbf{Mixed-block overlap.}  Strengthen Lemma~\ref{lem:block} by replacing the union
bound on mixed moduli with an overlap-sensitive bound.  The natural inputs are the CRT
dependency graph and exact block capacities $r(M)$.

\item \textbf{Higher-order BFF clusters.}  Generalize the triangle experiment to a certified
junction-tree or hypergraph inequality using all exactly known intersections of pairwise-disjoint
index sets.  The target is not merely a better numerical correction but a closed-form inequality
that remains monotone in the prime parameters.

\item \textbf{Independent certification.}  Reimplement the 105-case exponent certificate in a
second system and conduct a targeted literature search for the exact $218{,}243{,}025$ corollary
and maximum-spanning-forest formulation before making any novelty claim.
\end{enumerate}

\section{Reproducibility map}

\begin{center}
\begin{tabular}{lll}
\toprule
File & Role & Evidence level\\
\midrule
\texttt{bff\_core.py} & BFF formulas + optimized forest & published + derived\\
\texttt{verify\_218m\_bound.py} & 105 exact exponent cases & exact computation\\
\texttt{certificate\_218m.csv} & human-auditable case table & exact computation\\
\texttt{msf\_refinement.py} & optimized forest edges & derived + exact computation\\
\texttt{triangle\_exploration.py} & one-triangle cluster & derived/exploratory\\
\texttt{tower\_capacity.py} & block-capacity lemma/test & derived/exploratory\\
\texttt{fiber\_relaxation.py} & 3-adic capacity diagnostic & exploratory/exact\\
\texttt{research\_log.md} & provenance and next steps & audit trail\\
\bottomrule
\end{tabular}
\end{center}

\bibliographystyle{alpha}
\bibliography{references}

\end{document}
